\section{\gfaz as a Compute Engine}\label{sec:compute}

The previous sections present \gfaz as a compressor. The same representation also
supports a broader systems role: \gfaz can act as a compute substrate for
path-centric pangenome analyses. Many downstream tasks do not require arbitrary
mutation of the graph. Instead, they repeatedly scan haplotype traversals and
fold node visits into a compact accumulator: a coverage histogram, a
window-by-sample matrix, or an allele table relative to a reference path. Because
\texttt{.gfaz} already stores P/W traversals as grammar-compressed integer
streams, these analyses can consume the compressed traversal store directly and
avoid reconstructing the full GFA text or a whole-graph object.

This section describes the compute-engine design. The central abstraction is a
streaming traversal primitive that expands grammar rules on demand and emits
orientation-signed node IDs to an analysis-specific visitor
(Section~\ref{subsec:ce-model}). A bounded rule-leaf cache exposes a
memory/recompute trade-off for workloads that repeatedly touch large rules
(Section~\ref{subsec:ce-cache}). The resulting memory footprint is the compressed
traversal store plus the workload accumulator, rather than the decompressed
traversal graph (Section~\ref{subsec:ce-memory}). We instantiate the model in
three commands---\texttt{growth}, \texttt{pav}, and \texttt{deconstruct}---that
share the same decode substrate but use different accumulators
(Section~\ref{subsec:ce-instances}).

\subsection{Motivation: Avoiding Materialization}\label{subsec:ce-motivation}

Pangenome analysis tools often begin by materializing a representation larger
than the statistic they ultimately compute. Panacus~\cite{parmigiani2024panacus}
parses GFA path and walk records into item tables before deriving growth and core
curves. odgi~\cite{guarracino2022odgi} builds a succinct whole-graph object
before running operations such as \texttt{odgi pav}. \texttt{vg deconstruct}
loads a graph and snarl decomposition before emitting a reference-relative VCF
from traversals~\cite{garrison2018variation,hickey2020genotyping}. These
representations are appropriate for general graph operations, but for
path-centric folds the resident object can dominate both memory and time.

\gfaz targets a narrower but common execution pattern: compressed traversal
stream $\rightarrow$ small accumulator $\rightarrow$ biological statistic. The
design deliberately does not try to replace topology-heavy graph libraries.
Instead, it exploits the fact that the dominant GFA component---haplotype
traversals---is already encoded in a streamable form. The analysis pays only for
the compressed traversal columns, the rulebook needed to expand them, and the
state required by the specific query.

\subsection{Execution Model: Streaming Leaf Expansion}\label{subsec:ce-model}

\paragraph{Compressed traversal store.}
The \texttt{.gfaz} container stores P traversals and W traversals as two flat
\texttt{int32} arrays. Each record has an encoded-length column, and the original
pre-compression traversal lengths are retained for workflows that need a
materialized fallback. ZSTD decompression of these blocks yields
grammar-compressed integer streams containing literals and rule IDs, but not the
fully expanded node sequences. The engine slices the flat arrays into zero-copy
\texttt{HapSlice} views:
\[
  \{\texttt{pointer},\ \texttt{encoded length},\ \texttt{original length}\}.
\]
The grammar rulebook is stored as two parallel \texttt{int32} arrays,
\texttt{rules\_first} and \texttt{rules\_second}, delta-decoded once after ZSTD
decompression. The first rule ID, \texttt{min\_rule\_id}, is read from the layer
metadata; \texttt{max\_rule\_id} is derived from the rule count. A value whose
absolute ID lies in this half-open range is a rule, and all other values are
literals. This single-comparison classification is sound because compression
reserves the rule namespace strictly above every literal node ID
(Section~\ref{subsec:pipeline}): rule IDs are assigned sequentially across rounds
starting from \texttt{start\_id}, so a decoder never has to consult a side table
to tell a node from a rule.

\paragraph{Streaming primitive.}
The reusable operation is:
\[
  \textsc{StreamDecodedNodes}(\texttt{HapSlice},\ \textsc{Visit}).
\]
It emits the decoded orientation-signed node IDs of one traversal, in order, to a
caller-supplied visitor. For each encoded element $e$, the decoder performs one
boundary check. If $|e|$ is outside the rule range, $e$ is a literal. Otherwise
the decoder indexes the rule arrays at $|e|-\texttt{min\_rule\_id}$ and expands
the two children recursively. A negative rule reference denotes the
reverse-complement orientation: the children are visited in reverse order and
their signs are flipped. Thus the same rulebook serves both strands without
duplicating rules.

The inverse-delta transform is fused into streaming. For the common
\texttt{delta\_round=1} case, the decoder keeps a running prefix sum and calls
the visitor immediately after reconstructing each node ID. For
\texttt{delta\_round=0}, literals and rule leaves are already in node-ID space.
For uncommon higher delta counts, the implementation materializes one traversal,
applies the required inverse passes, and then visits the recovered nodes.

\begin{algorithm}[t]
  \caption{Streaming traversal expansion for \texttt{delta\_round} 0 or 1}
  \label{alg:ce-stream}
  \begin{algorithmic}[1]
    \REQUIRE slice $E$, rule range $[m,M)$, rule arrays $R_1,R_2$,
             delta round $D \in \{0,1\}$, visitor \textsc{Visit}
    \STATE $\textit{prev} \gets 0$
    \STATE \textbf{define} \textsc{Emit}($v$):
    \IF{$D = 1$}
      \STATE $\textit{prev} \gets \textit{prev} + v$; \textsc{Visit}(\textit{prev})
    \ELSE
      \STATE \textsc{Visit}($v$)
    \ENDIF
    \STATE \textbf{define} \textsc{Expand}($r$, \textit{reverse}):
    \STATE \quad $(a,b) \gets (R_1[r-m], R_2[r-m])$
    \STATE \quad $(c_1,c_2) \gets \textit{reverse}\,?\,(b,a):(a,b)$
           \COMMENT{reverse swaps child order}
    \FORALL{$c \in (c_1,c_2)$}
      \IF{$|c| \in [m,M)$}
        \STATE \textsc{Expand}($|c|$, $\,\textit{reverse}\oplus(c<0)$)
               \COMMENT{rule child: fold parent and child orientation}
      \ELSE
        \STATE \textsc{Emit}($\textit{reverse}\,?\,{-}c:c$)
               \COMMENT{literal child: flip sign under reverse}
      \ENDIF
    \ENDFOR
    \FORALL{$e \in E$ in order}
      \IF{$|e| \in [m,M)$}
        \STATE \textsc{Expand}($|e|$, $e < 0$)
      \ELSE
        \STATE \textsc{Emit}($e$)
      \ENDIF
    \ENDFOR
  \end{algorithmic}
\end{algorithm}

The cost of decoding one traversal is linear in the expanded traversal length.
Rule resolution is an array lookup, not a hash lookup, and recursion depth is
bounded by the number of grammar rounds. The visitor is templated in the
implementation, so the per-node fold can be inlined through the decode loop.

\subsection{Rule-Leaf Cache}\label{subsec:ce-cache}

On-demand recursion avoids materializing all expanded traversals, but a large
rule may be referenced many times. For workloads where repeated rule expansion is
noticeable, \gfaz provides a bounded \texttt{RuleLeafCache}. The cache stores the
fully expanded forward leaf sequence for any rule whose expansion fits within a
user-specified byte budget. Reverse expansion is derived by scanning the cached
leaf sequence backward and negating signs, so no reverse copy is stored.

Cache construction is bottom-up over rule IDs. A rule is admitted only if its
children have already been expanded or can be represented lazily and if the
projected leaf sequence fits within the remaining budget. Rules that do not fit
fall back to recursive expansion, preserving correctness for every budget. The
cache is built once, then shared read-only by all decoding threads. A budget of
zero recovers fully lazy expansion; larger budgets trade memory for less repeated
recursion. In the current implementation, \texttt{pav} and \texttt{deconstruct}
enable this cache, while \texttt{growth} uses pure recursive expansion.

\subsection{Memory Model}\label{subsec:ce-memory}

The compute engine replaces decompressed traversal materialization with four
resident components:
\begin{enumerate}[noitemsep, topsep=2pt, leftmargin=*]
  \item \textbf{Grammar-compressed traversal store}: ZSTD-decompressed but still
        grammar-encoded P/W integer streams, encoded-length columns, and the
        rule arrays.
  \item \textbf{Optional rule-leaf cache}: bounded by a configurable byte
        budget.
  \item \textbf{Analysis accumulator}: the state required by the specific query,
        such as a node coverage vector, a CSR node-to-group index, or a per-site
        allele table.
  \item \textbf{Sequence payload when needed}: \texttt{deconstruct} additionally
        keeps segment DNA resident so it can spell REF and ALT alleles.
\end{enumerate}

The key distinction is that the traversal store remains grammar-compressed.
Memory still scales with the size of the compressed traversal columns and with
the query accumulator, but not with a fully expanded list of all path and walk
steps. This property is most useful for workloads that scan many haplotypes but
retain only aggregate state.

\subsection{Parallelization}\label{subsec:ce-parallel}

Traversal decoding is parallel over \texttt{HapSlice} objects. The flat traversal
arrays, rulebook, segment sequence payload, and rule-leaf cache are read-only
after setup, so distinct threads can decode distinct traversals without
synchronizing on the decode path. Because expanded traversal lengths vary widely
across haplotypes, the parallel loops use dynamic OpenMP scheduling for load
balance. Synchronization is pushed into the workload-specific fold, and the three
commands cover distinct accumulator policies:
\begin{itemize}[noitemsep, topsep=2pt, leftmargin=*]
  \item \textbf{\texttt{growth}} parallelizes over haplotype groups and shares a
        single node coverage array. Each thread keeps a private last-seen stamp
        array to suppress duplicate counts within a group, and performs an atomic
        increment only on the first visit of a node in that group; the stamp is a
        rolling byte that is reset when it wraps, avoiding a per-group clear of
        the stamp array. The coverage-to-histogram reduction is then computed with
        per-thread local histograms merged in a short critical section, and the
        closed-form growth curve is evaluated in parallel over cohort sizes.
  \item \textbf{\texttt{pav}} decodes each traversal into a private sorted-unique
        node set (lock-free, per-slice slots), reference slices additionally
        capturing their ordered node stream. A serial prefix-sum builds the CSR
        node-to-group index, after which the per-contig reference sweep runs in
        parallel and updates the output matrix with atomics; these are
        contention-free in the common case because distinct contigs map to
        disjoint window-index ranges.
  \item \textbf{\texttt{deconstruct}} first fixes sites with a single
        linear-time topology pass (the biconnected decomposition below), which
        decodes no haplotype, then streams haplotypes in parallel into per-thread
        observation buffers concatenated after the parallel region; site assembly
        and VCF formatting are a subsequent serial pass over the per-site
        observation groups.
\end{itemize}

This organization separates the shared decoding substrate from the accumulator
policy, and isolates each command's serial stages (the CSR build for \texttt{pav},
site assembly for \texttt{deconstruct}). It also leaves a clear GPU path:
per-slice rule expansion is regular integer work over contiguous arrays, while the
fold can be specialized for each analysis.

\subsection{Three Instantiations}\label{subsec:ce-instances}

All three implemented commands share the same setup. They deserialize
\texttt{CompressedData}, ZSTD-decompress the rulebook and P/W traversal blocks,
delta-decode the rule arrays, build \texttt{HapSlice} views, and reconstruct only
the metadata columns needed for grouping or reference lookup.


\noindent\textbf{Growth.}
This command computes Panacus-style union growth for nodes. P/W records are
grouped by one of several identities: per record, PanSN sample,
sample-haplotype, or sample-haplotype-sequence. For P-lines the identity is parsed
from the path name with a PanSN matcher (mirroring Panacus, including the
stripping of trailing \texttt{:start-end} coordinate suffixes); for W-lines the
sample, haplotype index, and sequence ID are already stored as separate columns
and need no parsing. Streaming one group updates a coverage array
$\texttt{cov}[v]$, where each node is counted at most once per group. The coverage
vector is reduced to a histogram $\texttt{hist}[c]=|\{v:\texttt{cov}[v]=c\}|$. For
a pangenome with $N$ groups, the expected number of nodes observed in a random
subset of size $k$ is
\[
  \texttt{growth}[k]
    =
  \sum_c \texttt{hist}[c]
    \left(1-\frac{\binom{N-c}{k}}{\binom{N}{k}}\right).
\]
The ratio $\binom{N-c}{k}/\binom{N}{k}$ is evaluated as the stable product
$\prod_{i=0}^{c-1}(N-k-i)/(N-i)$ rather than by forming the binomial coefficients,
which would overflow at pangenome $N$. The accumulator is linear in the number of
nodes, and this fold matches Panacus union-growth semantics (count $=$ node,
coverage $\ge 1$, quorum $\ge 0$).

\noindent\textbf{PAV.}
This command computes a window-by-group matrix over BED intervals on reference
paths or walks, in three passes over a shared read-only decode state. \emph{Pass~1}
decodes every traversal once into a private sorted-unique node set; reference
traversals (those a BED \texttt{chrom} resolves to) additionally retain their
ordered node stream, captured as a side effect of the same decode so no reference
is decoded twice. \emph{Pass~2} unions the per-slice node sets of each group and
flattens the result into a compressed-sparse-row $\text{node}\rightarrow\text{groups}$
index: a counting pass tallies the groups per node, an exclusive prefix sum forms
the row offsets, and a cursor scatter fills the group-ID column. CSR is chosen over
a $\text{node}\rightarrow\text{list}$ map to remove the per-node container header
overhead, which is significant when the node count reaches hundreds of millions.
The rule-leaf cache is released after Pass~1, since Passes~2--3 perform no further
grammar expansion. \emph{Pass~3} sweeps each reference stream, maintaining a
running genomic offset from the segment-length column so that each visited node
maps to a half-open base interval; the BED ranges for that contig, pre-sorted by
start, are intersected against the node interval with an advancing range cursor,
and the resulting $(\text{window},\text{node})$ pairs are deduplicated. For each
unique pair, the node length is added to the window denominator and to the
numerator of every group in that node's CSR row. Groups are formed by the same
identity modes as \texttt{growth} (per record, sample, or sample-haplotype). A
node counts at most once per group, matching \texttt{odgi pav} semantics; the
result is node-length-weighted PAV ratios computed without a whole-graph object.

\noindent\textbf{Deconstruct.}
The \texttt{deconstruct} command is the most topology-sensitive instance of the
compute engine. It derives a reference-relative VCF from the compressed traversal
store without materializing all haplotype paths or constructing a general-purpose
whole-graph index. Each reference contig is decoded once to an orientation-signed
node stream \texttt{ref\_nodes}; a prefix sum over segment lengths maps reference
node positions to genomic offsets and gives the contig length used in the VCF
header. Site discovery and allele observation are separated. The paper-relevant
\texttt{--vg-compat} mode defines sites from top-level snarls with a new global
BCC algorithm over the topology stored in the compressed container. Allele
observation remains the same path-centric fold as the other compute-engine
commands: each compressed haplotype is streamed once through compact
site-boundary state.

Existing deconstruction pipelines typically combine graph loading, snarl
decomposition, and path inspection inside a materialized graph representation. In
\gfaz, the BCC algorithm keeps site discovery off the haplotype hot path. It
builds an orientation-aware node-end graph from the stored L-line topology: each
segment $s$ contributes two vertices, $s.5'$ and $s.3'$, joined by a \emph{black}
edge, and each GFA link contributes one \emph{grey} edge between the source
departure end and target arrival end. This is the bidirected ``interval'' graph
used by snarl theory, where snarls in a bidirected genome graph correspond to
biconnected blocks in a doubled/node-end representation~\cite{paten2018superbubbles}.
The black segment edges are essential for inversions: a single-node inversion
closes through the segment's two ends and is recovered as one biconnected block,
whereas a collapsed segment graph can expose the inverted segment as an
articulation point and split the site.

\gfaz\ runs an iterative Hopcroft--Tarjan decomposition over this node-end graph
(Algorithm~\ref{alg:ce-bcc}). The traversal maintains discovery times, low-link
values, and an edge stack; when returning from child $w$ to parent $u$, the
condition $\textit{low}(w)\ge\textit{disc}(u)$ pops one maximal biconnected block.
The pass is iterative rather than recursive to avoid call-stack limits on
chromosome-scale graphs, and runs in $O(V+E)$ time with linear auxiliary space.
Each non-trivial block is projected onto \texttt{ref\_nodes}: if the first and
last reference positions touched by the block define a clean span whose positions
all belong to the block and contain no repeated segment, the span is emitted as a
candidate top-level snarl. This filter rejects cyclic or ambiguous reference
traversals, while nested child bubbles and chains of leaf bubbles collapse into
their enclosing top-level snarl, matching the default record granularity of
\texttt{vg deconstruct}.

\begin{algorithm}[t]
  \caption{Top-level snarl extraction via biconnected decomposition}
  \label{alg:ce-bcc}
  \begin{algorithmic}[1]
    \REQUIRE links $L$; reference node stream \texttt{ref}; segment count $n$
    \STATE build node-end graph $G$: vertices $\{s.5',s.3'\}$ per segment $s$;
           \emph{black} edge $(s.5',s.3')$; \emph{grey} edge
           $(\textsc{dep}(a),\textsc{arr}(b))$ per link $a{\to}b$
    \STATE $t \gets 0$;\quad edge stack $\Sigma \gets [\,]$;\quad
           $\textit{disc}[\cdot]\gets 0$
    \STATE \textbf{procedure} \textsc{Dfs}($u$, $\textit{pe}$): \COMMENT{entered via edge \textit{pe}}
    \STATE \quad $\textit{disc}[u]\gets\textit{low}[u]\gets (t\gets t{+}1)$
    \STATE \quad \textbf{for each} incident edge $e=(u,w)\neq \textit{pe}$:
    \STATE \qquad \textbf{if} $\textit{disc}[w]=0$: \quad push $e$ on $\Sigma$;\;
           \textsc{Dfs}($w$,$e$);\;
           $\textit{low}[u]\gets\min(\textit{low}[u],\textit{low}[w])$
    \STATE \qquad\quad \textbf{if} $\textit{low}[w]\ge\textit{disc}[u]$:
           pop $\Sigma$ down to and including $e$ as block $B$;\; \textsc{Project}($B$)
    \STATE \qquad \textbf{else if} $\textit{disc}[w]<\textit{disc}[u]$:
           push $e$ on $\Sigma$;\; $\textit{low}[u]\gets\min(\textit{low}[u],\textit{disc}[w])$
    \STATE \textbf{for each} reference segment $s$ and end $x\in\{s.5',s.3'\}$ with $\textit{disc}[x]=0$: \textsc{Dfs}($x$, \textsc{nil})
    \STATE \textbf{procedure} \textsc{Project}($B$):
    \STATE \quad $S \gets$ segments incident to $B$;\quad
           $I \gets \{\,i : \texttt{ref}[i]\in S\,\}$
    \STATE \quad \textbf{if} $|B|\ge 2$ \textbf{and} $|I|\ge 2$:
           $(\ell,h)\gets(\min I,\max I)$
    \STATE \qquad \textbf{if} \texttt{ref}$[\ell..h]$ repeats no segment \textbf{and}
           $\texttt{ref}[i]\in S\ \forall i\in[\ell,h]$:
           emit snarl $(\texttt{ref}[\ell],\texttt{ref}[h])$
    \STATE reduce emitted snarls to a non-overlapping chain by position
  \end{algorithmic}
\end{algorithm}

Once BCC sites are fixed, \gfaz\ returns to the compressed traversal substrate.
For every top-level snarl, it stores forward and reverse entrance maps. The
forward entrance maps the reference source boundary to the snarl and expects the
reference sink as its exit; the reverse entrance maps the reverse-complement sink
to the same snarl and expects the reverse-complement source, allowing inverted
haplotype traversals to be normalized back into the reference frame. During
streaming, a small state machine is either closed or associated with one open
snarl (Algorithm~\ref{alg:ce-snarl}). A boundary opens a site, ordinary decoded
nodes are appended to a small interior buffer, and the matching exit emits that
interior tagged by the haplotype slice. If traversal was reversed through the
site, the buffer is scanned backward and signs are negated before emission; a
boundary may close one site and immediately open the next, which handles adjacent
snarls without retaining the surrounding path.

\begin{algorithm}[t]
  \caption{Streaming allele observation for one haplotype (snarl mode)}
  \label{alg:ce-snarl}
  \begin{algorithmic}[1]
    \REQUIRE slice $E$; per-snarl boundary maps \textsc{Entrance} (node-side
    $\rightarrow$ snarl, forward or reversed) and exit node-sides
    \STATE $open \gets \bot$; $\textit{interior} \gets [\,]$
    \STATE \textbf{for each} decoded node $u$ of $E$ (Algorithm~\ref{alg:ce-stream}):
    \STATE \quad $v \gets$ node-side vertex of $u$
    \IF{$open = \bot$}
      \STATE \textsc{TryOpen}($v$): if $v$ is a (forward/reversed) entrance, set
      $open$, \textit{reversed}, $\textit{exit}$, $\textit{interior}\gets[\,]$
    \ELSIF{$v = \textit{exit}$}
      \STATE emit interior for snarl $open$ (reverse + negate if \textit{reversed}),
      tagged with this slice
      \STATE $open \gets \bot$; \textsc{TryOpen}($v$) \COMMENT{exit may open a chained snarl}
    \ELSE
      \STATE append $u$ to \textit{interior}
    \ENDIF
  \end{algorithmic}
\end{algorithm}

Per site, the reference interior spells allele $0$; each observed interior is
spelled by concatenating segment DNA (reverse-oriented nodes contributing reverse
complements) and deduplicated to an allele index. A site with no alternate allele
is dropped. Equal-length alleles are emitted as substitutions; otherwise the
record is left-anchored with the last base of the source anchor, following the
standard VCF indel convention. \texttt{INFO} carries \texttt{AC}/\texttt{AN}/
\texttt{AF} and the number of samples with data, and an optional
\texttt{--max-site-length} guard collapses very large spans to a symbolic
\texttt{<CPX>} record. Haplotypes collapse into VCF sample columns by the chosen
grouping mode, with one genotype slot per haplotype index so reference-tiling
contigs resolve to the sample's true ploidy rather than inflating it. Records are
emitted in sorted genomic order per contig, yielding a GFAz-native
reference-relative deconstruction workflow computed directly on compressed
traversals.

These three commands delimit the intended scope of the compute engine: analytics
whose hot path is traversal streaming plus a compact fold. Operations that require
arbitrary node-to-step queries, graph editing, or full topology transformation
remain better served by materialized graph toolkits.
